\appendix
\section{口径与数据说明}

\subsection{统计口径纪律}

\begin{itemize}
  \item 每个数字标注规模 $n$ 与样本框；不同 $n$/口径的读数不连成同一演化线（图注注明口径切换点）。
  \item 计分卡读数一律按获取样本集过滤后计算——stagerun 落盘记录为追加式，跨轮重跑会残留过期行，原始行计数会把未参与样本的跳过计入分母。
  \item 编译终态三级：纯净 = 无告警出 PDF；带瑕 = 带可接受瑕疵出 PDF（\texttt{acceptable}、\texttt{best\_effort}、\texttt{dirty} 三档）；联合口径出 PDF = 纯净 + 带瑕。
  \item 延迟为网关端到端秒；批量档入闸排队约 120s，拥塞期 p95 会被队列时间放大，读数与空闲档不可比。
  \item qualbench 评分由 LLM 评审（swe-2-max）产出，争议判例经第二评审（swe-2-high）仲裁；评审自身偏差未做人工标定，分数应读作相对分布而非绝对正确率。
\end{itemize}

\subsection{随档清单与原始产物索引}

本报告图表全部由随档 \texttt{data/} 原始数据生成，不依赖本目录以外的图片路径；但 \texttt{refs/} 仅随档上一期复盘的两张 SVG，引用源文档的时点快照未随迁——对应档案位置见下表与 README.md（「随档 \texttt{refs/} 全收录」的原始表述以 README 为准）。

\begin{center}\small
\begin{tabular}{@{}l>{\raggedright\arraybackslash}p{10.6cm}@{}}
\toprule
目录 & 内容 \\
\midrule
\texttt{data/} & 六条时间线 CSV（计分卡 / 解析 / 语料 / 规则测试 / 翻译质量 / 每日 commit）；\texttt{arms\_history.csv} 377 行评测流水（含各轮 run 目录名、样本量、头条读数）；\texttt{commits.json} 1{,}648 commit 全录；\texttt{milestones.md} D0--D5 大事记 \\
\midrule
\texttt{refs/} & 随档仅 \texttt{metrics-timeline.svg} + \texttt{assets-growth.svg}（上一期复盘两图）。其余源文档快照未随迁，对应档案：评测器规格（七臂定义/门槛）→ \texttt{docs/spec/benchmark.md}；语料规格（层构造/治理）→ \texttt{docs/spec/corpus.md}；逐库横评协议 → \texttt{docs/dev/bench-harness.md}；验证分层契约 L0--L3 → \texttt{bench/TIERS.md}；层语义/EVAL\_ONLY → \texttt{bench/corpus/MANIFEST*.md} 族；机制台账快照 → \texttt{bench/corpus/mechanisms.jsonl}；8 库横评（D0 定案证据）→ \texttt{bench/archive-2026-09-20/results/00-grand-comparison.md}；本批 Markdown 总表 → \texttt{docs/research/methods/2026-09-19-bench-metrics-corpusv3.md}；上一期复盘正文与车队台账快照 → 迁移前存档 \texttt{tmp/old-docs-2026-09-20/docs/research/metrics-2026-09-19/refs/} \\
\bottomrule
\end{tabular}
\end{center}

各臂原始评测产物（编译工作区、译文 jsonl、PDF 中间件）为大体量文件不入档；溯源时按 \texttt{data/arms\_history.csv} 的 run 目录名在评测产物存档区查找对应轮次。

\subsection{过程事故记录（评测工具链教训）}

本批测试记录的评测驱动层教训，供后续批次回避：xlatbench 的 \texttt{manifest.parent} 必须等于语料根（跨目录清单会致格解析失败）；\texttt{--per-kind} 是全局限额非分层配额；pgrep 监控须用 \texttt{[x]pattern} 括号形防自匹配；compilebench \texttt{--gen-sample} 为分段设计（采样后即返回，编译需无参二次调用）；e2e 同 \texttt{--tag} 冻结样本，重抽样须换 tag+seed；gullet \texttt{--out} 续跑会保留过期错误行、须换新目录；stagerun 下游阶段须显式传 \texttt{--ids} 防跨层欠采；多会话共仓提交须走私有 \texttt{GIT\_INDEX\_FILE}（共享索引有 sweep/reset 竞争窗）。
